\section{A further simultaneous center-product test}
\label{sec:center-product-followup}

The following small exact continuation changes no general exclusion or
realization count.  It strengthens the necessary center data for
\texttt{R8\_3\_721} and \texttt{R8\_3\_1154}.  A surviving tuple below
is not a center fusion ring or a categorical model.  The previously
verified signed-sphericalization and Galois-pseudo-unitarity gates for
these two rings make the pseudo-unitary tests necessary for general
complex categorification.

\begin{lemma}[Simultaneous pair-product Gram condition]
\label{lem:followup-pair-gram}
Let $\mathcal C$ be pseudo-unitary, put $D=\operatorname{FPdim}(\mathcal C)$,
and let $A_0=\mathbf1,A_1,\ldots,A_s$ be self-dual simple objects of its
center whose complete unnormalized $S$ rows are known.  If $d_Z$ is the
dimension of a central simple $Z$, then
\[
 K_{ij,kl}=\frac1{D^2}\sum_Z
 \frac{S_{A_i,Z}S_{A_j,Z}S_{A_k,Z}S_{A_l,Z}}{d_Z^2}
 =\sum_Y N_{A_i,A_j}^{Y}N_{A_k,A_l}^{Y}.
\]
Subtract the summands for $Y=A_0,\ldots,A_s$ to obtain $G$.
Let $M$ have the restriction vectors of all individual remaining central
simples as columns, and let $T$ have the residual forgetful pair products
as columns.  Then their unknown pair multiplicities form a matrix
\[
             X\in\mathbb Z_{\ge0}^{r\times\binom{s+1}{2}},
             \qquad MX=T,\qquad X^{\mathsf t}X=G.
\]
Repeated restriction vectors are repeated columns of $M$; they are not
identified as central simples.
\end{lemma}

\begin{proof}
Apply the ordinary Verlinde formula twice and use
$S^2=D^2C$, where $C$ is charge conjugation.  Self-duality of the four
displayed objects gives the stated bilinear pairing.  No assumption
of a $C^*$ structure is used.  Subtracting the known simple coefficients
gives the Gram identity, and applying the tensor forgetful functor to
each pair gives $MX=T$.
\end{proof}

The normalization uses the center dimension $D^2$.  In the
noncommutative source case the complete distinguished row is
$S_{A_E,Z}=D\chi_E(FZ)/(f_E\dim E)$; the matrix-block constituent is
counted once as a central simple.  As an independent control, the
electric objects of $Z(\operatorname{Vec}(S_3))$ reproduce all $27$
ordered cubic coefficients and all $81$ ordered fourth moments of
$\operatorname{Rep}(S_3)$.  In particular the source block with
$\dim E=2$ and $f_E=3$ has normalization $6/(3\cdot2)=1$.
The categorical input is ordinary Verlinde, equivalently the regular
module specialization of Deshpande, \emph{On twisted Verlinde formulae
for modular categories}, Theorem~1.1, arXiv:1811.08447.

\subsection{Forced pair products for \texttt{R8\_3\_721}}

All five distinguished objects are self-dual: their complete source
trace characters satisfy $\chi_E(x^*)=\chi_E(x)$, so the corresponding
central $S$ rows are invariant under duality, and nondegeneracy identifies
each object with its dual.  Start with the four complete integral Gram
patterns already obtained in the frozen calculation, and retain both
symmetric character assignments.  All $35$ fourth moments for the four
nonunit low objects are nonnegative integers in each of these eight cases.

Order pairs as $11,12,13,14,22,23,24,33,34,44$.  In every case
the residual $G$ has rank four and $M$ has rank three.  If $X_0$ is the
exact minimum-norm solution of $MX=T$, then
\[
 G-X_0^{\mathsf t}X_0=\frac23h^{\mathsf t}h,
 \qquad h=(0,0,0,0,1,1,-1,1,-1,1).
\]
Thus every solution is $X=X_0+t h$, where $Mt=0$ and
$\sum t_i^2=2/3$.  Each $t_i$ lies in $(-1,1)$; integrality of one
coordinate fixes it modulo $\mathbb Z$.  Enforcing integrality and
nonnegativity in all ten coordinates gives a unique minimum-square
choice in each row.  The sum of those minima is already $2/3$.
Consequently the entire integral matrix $X$ is unique.  Direct exact
substitution checks both original identities.

Here are its row types, with the high-column indices of the frozen
Gram enumeration:
\[
\begin{array}{c|rrrrrrrrrr}
\text{type}&11&12&13&14&22&23&24&33&34&44\\\hline
0&2&2&2&1&2&3&1&2&1&1\\
1&4&4&4&2&4&6&2&4&2&2\\
3&2&2&2&1&3&2&0&3&0&0\\
4&4&4&4&2&4&4&2&4&2&0\\
5&4&4&4&2&4&4&2&4&2&1\\
8&8&8&8&4&8&8&4&8&4&0
\end{array}
\]
Denote the type-$0$-or-$1$ group by $\mathsf A$, the type-$3$
singleton by $\mathsf B$, the type-$4$-or-$8$ group by $\mathsf C$,
and the eight type-$5$ simples by $\mathsf D$.  The sizes
$m_A,m_C$ are independently one or four: a singleton $\mathsf A$
uses type~1 and four copies use type~0; a singleton $\mathsf C$ uses
type~8 and four copies use type~4.  The resulting center ranks are
$16,19,19,22$.

\subsection{Complete low-action extension of these patterns}

Write $N_i$ for fusion by $A_i$.  The known vectors
$e_0,\ldots,e_4$ and $N_iN_je_0$, $1\le i\le j\le4$, span a
nine-dimensional space.  Put $u=33+8\sqrt{17}$ and
$u'=33-8\sqrt{17}$.  The complete simultaneous eigenvalue tuples of
$(N_1,N_2,N_3,N_4)$ are
\[
\begin{array}{c|c}
\text{tuple}&\text{multiplicity}\\\hline
(u,u+1,u+1,(u+1)/2)&1\\
(u',u'+1,u'+1,(u'+1)/2)&1\\
(1,-2,2,1),\ (1,2,-2,1),\ (1,2,2,5)&1\text{ each}\\
(1,-2,-2,1)&m_A\\
(1,2,2,-3)&1\\
(1,0,0,-1)&m_C\\
(-1,0,0,0)&8
\end{array}
\]
Every tuple has positive vacuum weight $\sum d_Z^2/D^2$.
The vacuum cyclic space therefore has dimension nine, so it is exactly
the already known span.  Select its basis columns $V$, remembering
the corresponding words $p_\alpha$ of degree at most two.  With
$G_0=V^{\mathsf t}V$, its forced action, extended by zero on the
orthogonal complement, is
\[
 F_i=VG_0^{-1}B_iG_0^{-1}V^{\mathsf t},\qquad
 (B_i)_{\alpha\beta}=D^{-2}\sum_Zd_Z^2
 p_\alpha(\lambda_Z)\lambda_{i,Z}p_\beta(\lambda_Z).
\]
These degree-at-most-five moments give exact rational matrices.
Let $Q=I-VG_0^{-1}V^{\mathsf t}$.  Its range is the sum of the
within-group difference spaces of $\mathsf A,\mathsf C,\mathsf D$.
The remaining simultaneous spectral projectors $P_A,P_C,P_D$ have
ranks $m_A-1,m_C-1,7$, are mutually orthogonal, and sum to $Q$.
Every possible extension is consequently
\begin{equation}\label{eq:followup-hidden-actions}
\begin{split}
N_1&=F_1-Q+2(P_A+P_C),&N_2&=F_2-2P_A,\\
N_3&=F_3-2P_A,&N_4&=F_4+P_A-P_C.
\end{split}
\end{equation}
This describes the entire remaining spectral domain, not a choice of
canonical projectors.

\begin{proposition}\label{prop:r721-action-extension}
The patterns $(m_A,m_C)=(1,4)$ and $(4,1)$ are impossible.
The pattern $(1,1)$ has a unique nonnegative integral tuple of the four
necessary low fusion actions.  The pattern $(4,4)$ has exactly $46$
labelled such tuples, forming three classes under independent
permutations within the equal-restriction groups.
All surviving tuples satisfy the exact forgetful intertwining relations.
These assertions do not construct a full center fusion ring or a category.
\end{proposition}

\begin{proof}
For $(1,4)$, $P_A=0$, and the $\mathsf C$ block of the forced matrix
$N_2=F_2$ is $(9/2)J_4$, contradicting integrality.
For $(4,1)$, $P_C=0$, and the fixed $\mathsf A$ block of
$N_1+N_2=F_1+F_2-Q$ is $(5/2)J_4-I_4$, again contradicting
integrality.  Both statements hold for both symmetric character
assignments.  For $(1,1)$ both projectors vanish; exact substitution
in \eqref{eq:followup-hidden-actions} gives the asserted unique tuple.

It remains to give a complete finite analysis when both sizes are four.
The fixed within-group blocks of $F_1,F_2,F_3,F_4$ are respectively
\[
 \mathsf A:\quad (5/4)J_4,\ J_4,\ J_4,\ (3/4)J_4;
 \qquad
 \mathsf C:\quad (17/4)J_4,\ (9/2)J_4,\ (9/2)J_4,\ (9/4)J_4.
\]
Use the elementary identity $E_{ii}=\sum_jE_{ij}^2$ for an
orthogonal projector $E$.  Since $0\le P_A,P_C,P\le Q$ as
orthogonal subprojectors, their diagonal entries on $\mathsf A$ and
$\mathsf C$ are at most $Q_{ii}=3/4$.
All hidden projectors annihilate the
constant vector of each high group.

Set $P=P_A+P_C$.  Integrality of $N_1$ forces its diagonal entries
in either group to be $1/4$ or $3/4$, and its within-group
off-diagonal entries to be odd multiples of $1/4$.  In either case
the square of the diagonal plus three times $(1/4)^2$ already equals
the diagonal.  Hence every within-group off-diagonal absolute value
is $1/4$, and all entries outside that group vanish.

Nonnegative integrality of $N_2$ gives $(P_A)_{ii}\in\{0,1/2\}$
on $\mathsf A$, with half-integral within-group entries.  A zero
diagonal gives a zero row.  A diagonal $1/2$ and zero group row sum
force an off-diagonal entry at most $-1/2$; its norm identity is then
exhausted.  Thus $P_A$ on $\mathsf A$ is a sum $Q_A$ of disjoint
matching projectors $(e_j-e_k)(e_j-e_k)^{\mathsf t}/2$, with no
cross entries.  On $\mathsf C$, integrality gives diagonal $1/4$
or $3/4$ and odd-quarter off-diagonal entries.  The same norm argument
gives off-diagonal absolute values $1/4$ and no cross entries.

The diagonal integrality of $N_4=F_4+2P_A-P$ now forces every
diagonal of $P$ in both groups to be $3/4$.  Zero row sums force
all off-diagonals to be $-1/4$.  Therefore $P=E_A+E_C$, the two
canonical difference projectors; their ranks exhaust its rank six,
so it vanishes on $\mathsf D$.  On $\mathsf C$ the complement
$E_C-P_A$ has diagonal zero or $1/2$ and off-diagonal zero or
$-1/2$, so it too is a matching projector, say $Q_C$.
The rank-three condition on $P_A$ says the two matchings have equal
size.  Exhaustively,
\[
 P_A=Q_A+E_C-Q_C,\qquad
 P_C=E_A-Q_A+Q_C,\qquad P_D=E_D.
\]
There are $1,6,3$ matchings on four labelled vertices of sizes
$0,1,2$, giving $1+6^2+3^2=46$ choices and three permutation classes.
All these matrices pass direct exact substitution in
\eqref{eq:followup-hidden-actions}.

For every surviving tuple the verifier checks symmetry, nonnegative
integrality, the unit column, all forced pair columns, pairwise
commutation, projector ranks and orthogonality, spectral trace and
square trace, and
$M_{\mathrm{full}}N_i=L_{F(A_i)}M_{\mathrm{full}}$.
This proves the stated necessary-action result.
\end{proof}

The high-block traces and squared Frobenius norms provide compact checks:
\[
\begin{array}{c|c|c}
\text{center rank}&(\operatorname{tr}H_i)_{i=1}^4&
 (\operatorname{tr}H_i^2)_{i=1}^4\\\hline
16&(48,56,56,28)&(2798,2880,2880,746)\\
22&(54,50,50,28)&(2804,2892,2892,752)
\end{array}
\]
The full matrices, source bindings, matching domains and exact replays are
retained with \texttt{quartic\_center\_challenge.py} and
\texttt{r721\_center\_extension.py}.  Checks remain active under Python
optimization.  The two surviving Gram patterns still require products
of arbitrary high objects, a full based ring and categorical coherence.
